% Theorem C': the OT leak toward a prior with a different number of atoms (PARADIGM.md, section 4b).
\section{Theorem C$'$: the OT leak toward a prior with a different number of atoms}\label{sec:thmCprime}

Theorem C needs a prior with as many atoms as the network. Here the prior is any $\nu\in\Prob_2(\ThKe)$, for example
the empirical measure of a network of a different size (the transfer problem of Section~\ref{sec:program}, branch 8),
and the leak points at the barycentre of the prior mass that the optimal plan assigns to each particle. The
dissipation toward the prior survives up to the quantisation error, the fixed points of the pure leak are the
equal-mass centroidal configurations of $\nu$ (the leak is Lloyd's algorithm in continuous time), and the Lyapunov
bias is measured by the distance of the ideal from the \emph{centroids} of its prior cells, not by $W(\mu^*,\nu)$.

\paragraph{The law.} Let $\nu\in\Prob_2(\ThKe)$, for instance $\nu = \frac1{m_0}\sum_{j=1}^{m_0}\delta_{\vartheta^0_j}$
with any $m_0\ge1$. For $\mu = \frac1m\sum\delta_{\vartheta_i}$, let $\gamma$ be a vertex optimal plan from $\mu$ to
$\nu$ for the cost $\norm{\cdot}^2_{\Ginv}$ (for atomic $\nu$ a vertex of the transportation polytope; for general
$\nu$ any optimal plan), and let
\[
  \bar T_i := m\int\vartheta'\dd\gamma(\vartheta_i,\vartheta')
\]
(the barycentric projection of $\gamma$ at $\vartheta_i$; for atomic $\nu$,
$\bar T_i = m\sum_j\gamma_{ij}\vartheta^0_j$),
\begin{gather*}
  V_{\mathrm{bar}}^2 := \frac1m\sum_i\norm{\vartheta_i - \bar T_i}^2_{\Ginv},\qquad
  \mathrm{Var}_\gamma := \int\norm{\vartheta' - \bar T(\vartheta)}^2_{\Ginv}\dd\gamma(\vartheta,\vartheta'),\\
  \text{so that}\qquad W(\mu,\nu)^2 = V_{\mathrm{bar}}^2 + \mathrm{Var}_\gamma.
\end{gather*}
The \textbf{barycentric OT leak} is
\begin{equation}\label{eq:thmCprime:Lsigbar}
  \dot\vartheta_i = \Proj\bigl(\vartheta_i,\ G\nabla g_t(\vartheta_i) - \sigma_m(\vartheta_i - \bar T_i)\bigr).
  \tag{$\mathrm{L}_{\sigma,\mathrm{bar}}$}
\end{equation}
For atomic $\nu$ the plan is a vertex of the transportation polytope (row sums $1/m$, column sums $1/m_0$), computed
by linear programming (\code{transport.optimal_plan}, \code{transport.barycentric_projection}); for $m_0 = m$ it is a
permutation and \textup{(L$_{\sigma,\mathrm{bar}}$)} is $\Lsig$. The right-hand side is discontinuous where the
optimal plan is not unique, and here, unlike Section~\ref{sec:thmC} (Lemma C.3), the plan does switch along solutions
(numerics below), so two solution concepts are used. \emph{Vertex-selection solutions} are absolutely continuous
$\vartheta(t)$ with
$\dot\vartheta_i = \Proj\bigl(\vartheta_i,\ G\nabla g_t - \sigma_m(\vartheta_i - \bar T_i^{\gamma(t)})\bigr)$ for a.e.\
$t$ and $\gamma(t)$ a vertex optimal plan at $\vartheta(t)$; the time-stepped implementation produces these (switches
at step times), and their existence for all time is assumed. \emph{Generalised solutions} allow $\gamma(t)$ to be any
measurable selection of optimal plans. The optimal face $F(\vartheta)$ of the transportation LP is a compact convex
polytope, upper semicontinuous in $\vartheta$ (finitely many vertices, continuous costs), and
$\gamma\mapsto\bar T^\gamma$ is linear, so the set of generalised velocities is compact, convex and upper
semicontinuous in $\vartheta$: it contains the Filippov regularisation of the vertex law, and generalised solutions
exist on $[0,\infty)$ by the standard theorem for such differential inclusions (the projection is handled as in
Section~\ref{sec:thmC}). Every statement below holds for both concepts, with $\bar T$ and $V_{\mathrm{bar}}$ computed
from the plan the solution uses. Since $\ThKe$ is convex, $\bar T_i\in\ThKe$, so property (P) applies with
$\vartheta' = \bar T_i$.

\paragraph{Notation.} $\eta := W(\mu_t,\mu^*_t)$, $b := \sup_t W(\mu^*_t,\nu)$, and
\[
  b_{\mathrm{bar}}(t)^2 := \frac1m\sum_j\norm{\vartheta^*_j(t) - \bar T^*_j(t)}^2_{\Ginv}
  = W(\mu^*_t,\nu)^2 - \mathrm{Var}_{\gamma^*(t)},\qquad
  b_{\mathrm{bar}} := \sup_t b_{\mathrm{bar}}(t),
\]
the \emph{centroid defect of the ideal}, computed with an optimal plan $\gamma^*(t)$ from $\mu^*_t$ to $\nu$.
$\Delta_{\mathrm{gap}}(t)$ is the gap between the optimal and the next vertex cost of the transportation polytope at
$\mu^*_t$ ($\min\emptyset := +\infty$), and $\Delta_{\mathrm{gap}} := \inf_t\Delta_{\mathrm{gap}}(t)$.

\begin{thm}{Theorem}{C$'$}
Under \textup{(S)}, along every vertex-selection or generalised solution of \textup{(L$_{\sigma,\mathrm{bar}}$)}:
\begin{enumerate}[label=\textup{(\roman*)}]
  \item \emph{Dissipation toward the centroidal configurations.} With $Z := \frac12W(\mu_t,\nu)^2$ and $\Gamma$ as in
    Theorem C(i),
    \[
      D^+Z \le -\sigma_mV_{\mathrm{bar}}^2 + V_{\mathrm{bar}}\Gamma
      \le -2\sigma_m\bigl(Z - \tfrac12\mathrm{Var}_\gamma\bigr) + (2Z)^{1/2}\Gamma
      \qquad\text{for a.e.\ } t.
    \]
    For the pure leak \textup{(}$\Gamma\equiv0$\textup{)}, $W(\mu_t,\nu)$ is non-increasing, strictly decreasing
    wherever the plan in use is not centroidal \textup{(}$V_{\mathrm{bar}}>0$\textup{)}, and at every $\omega$-limit
    point $\vartheta_\infty$ of a solution on the compact $\ThK$ some optimal plan $\gamma_\infty$ at
    $\vartheta_\infty$ is centroidal: $\vartheta_{\infty,i} = \bar T_i^{\gamma_\infty}$ for all $i$ --- a vertex plan
    for a vertex-selection solution. Such configurations are the fixed points of the equal-mass Lloyd iteration
    $\vartheta\leftarrow\bar T^\gamma(\vartheta)$ with an optimal \textup{(}vertex\textup{)} plan; where the optimal
    plan is unique they are the critical points of the quantisation energy
    $\vartheta\mapsto\frac12W^2\bigl(\frac1m\sum\delta_{\vartheta_i},\nu\bigr)$, which is differentiable there, and the
    pure leak is its $G$-gradient flow with the factor $m$:
    $\dot\vartheta_i = -\sigma_m m\,G\,\partial_{\vartheta_i}\bigl[\tfrac12W^2\bigr]$. They include degenerate
    configurations with coincident particles. For $m_0 = m$, $\mathrm{Var}_\gamma = 0$ and \textup{(i)} is
    Theorem C(i).
  \item \emph{Induced coupling.} Let $\pi^*$ be optimal from $\mu$ to $\mu^*$, and let $\gamma$ be the plan the law
    uses at $t$. Put $C_{\mathrm{bar}} := \frac1m\sum_i\norm{\vartheta^*_{\pi^*(i)} - \bar T_i}^2_{\Ginv}$. Then
    \[
      \sqrt{C_{\mathrm{bar}}} \le b(t) + 2\eta \quad\text{always},
    \]
    and, for atomic $\nu$, $C_{\mathrm{bar}} = b_{\mathrm{bar}}(t)^2$ whenever
    $4\eta(b(t)+\eta) < \Delta_{\mathrm{gap}}(t)$ and $\gamma^*(t)$ is unique.
  \item \emph{Lyapunov inequality.} With $C_{\mathrm{bar}}$ as in \textup{(ii)},
    \[
      D^+\VOT \le -k_e\norm{e}^2 + \norm{e}\bigl(\varepsilon_{\max} + K\lambda_G\eta^2\bigr) + V_*\eta
      - \sigma_m\eta^2 + \sigma_m\eta\sqrt{C_{\mathrm{bar}}},
    \]
    and under the gap condition of \textup{(ii)} the leak terms are at most
    $-(\sigma_m/2)\eta^2 + (\sigma_m/2)b_{\mathrm{bar}}^2$.
  \item \emph{Local exponential estimate.} For atomic $\nu$, Theorem C(iv) holds verbatim with $b$ replaced by
    $b_{\mathrm{bar}}$: under $4\bar\eta(b+\bar\eta) < \Delta_{\mathrm{gap}}$, uniqueness of $\gamma^*(t)$, and
    $8K^2\lambda_G^2\bar\eta^2 \le k_e\sigma_m$, with $\alpha := \min(k_e,\sigma_m/2)$ and
    $\beta := \varepsilon_{\max}^2/k_e + (\sigma_m/2)b_{\mathrm{bar}}^2 + 2V_*^2/\sigma_m$, if
    $\max\{\VOT(0),\beta/\alpha\} < \bar\eta^2/2$ then $\VOT(t) \le e^{-\alpha t}\VOT(0) + \beta/\alpha$ and
    $\{\VOT\le c\}$ is forward invariant for $c\in\bigl[\max\{\VOT(0),\beta/\alpha\},\ \bar\eta^2/2\bigr)$.
  \item \emph{What the bias measures.} $b_{\mathrm{bar}} \le b$, with $b_{\mathrm{bar}} = 0$ exactly when the ideal is
    itself an equal-mass centroidal configuration of $\nu$. So a prior with many atoms whose centroidal
    $m$-quantisation coincides with the ideal anchors the network without bias, although $W(\mu^*,\nu) > 0$; a prior
    with the right number of atoms but the wrong positions has $b_{\mathrm{bar}} = b$.
\end{enumerate}
\end{thm}
\status{Proved and checked: dissipation at every step and convergence of the pure leak to an equal-mass Lloyd fixed
point, the coupling bound and plan stability, the Lyapunov inequality along the closed loop
(\code{tests/test_barycentric.py}, \code{exp_barycentric_leak.py}). Novelty (literature audit): new as a control
construction; the quantisation side is Lloyd/CVT theory; it reduces to Theorem C for $m_0=m$.}

\paragraph{Scope of (ii)--(iv).} The gap condition $4\eta(b+\eta) < \Delta_{\mathrm{gap}}(t)$ asks the network to be
closer to the ideal than the optimal vertex of the LP is to the next one. For a prior with many atoms,
$\Delta_{\mathrm{gap}}$ is at most the cost of exchanging two atoms between two cells, which is $O(1/m_0)$, so the
condition holds only in a small neighbourhood of the ideal; in the closed-loop experiment below $4\eta(0)(b+\eta(0))$
exceeds that bound by factors 83--175, and (iv) applies to no tested instance (the table reports the exchange bound on
$\Delta_{\mathrm{gap}}$, $\eta(0)$ and $4\eta(0)(b+\eta(0))$ for each run). What the experiment confirms is the
fixed-point structure, (i) and (v): the loop settles at the centroidal configuration of its prior, whose distance from
the ideal is $b_{\mathrm{bar}}$.

\begin{proof}
\textit{(i)} At a.e.\ $t$ the law uses an optimal plan $\gamma$, and for $h>0$ the same plan couples $\mu_{t+h}$ and
$\nu$, so $W(\mu_{t+h},\nu)^2 \le \sum_{ij}\gamma_{ij}\norm{\vartheta_i(t+h) - \vartheta^0_j}^2_{\Ginv}$ (for general
$\nu$, $\int\norm{\vartheta_i(t+h) - \vartheta'}^2\dd\gamma$). Differentiating the right-hand side at $h = 0$,
\[
  D^+Z \le \sum_{ij}\gamma_{ij}\ip{\vartheta_i - \vartheta^0_j}{\Ginv\dot\vartheta_i}
  = \frac1m\sum_i\ip{\vartheta_i - \bar T_i}{\Ginv\dot\vartheta_i},
\]
because $\sum_j\gamma_{ij}(\vartheta_i - \vartheta^0_j) = \frac1m(\vartheta_i - \bar T_i)$. Apply (P) with
$\vartheta' = \bar T_i\in\ThKe$ and $y = \nabla g_t - \sigma_m\Ginv(\vartheta_i - \bar T_i)$:
$\ip{\vartheta_i - \bar T_i}{\Ginv\dot\vartheta_i} \le \ip{\vartheta_i - \bar T_i}{\nabla g_t(\vartheta_i)}
- \sigma_m\norm{\vartheta_i - \bar T_i}^2_{\Ginv}$. Average and use Cauchy--Schwarz for the first term:
$\frac1m\sum\ip{\vartheta_i - \bar T_i}{\nabla g_t} \le V_{\mathrm{bar}}\Gamma \le (2Z)^{1/2}\Gamma$. The Pythagorean
identity $W^2 = V_{\mathrm{bar}}^2 + \mathrm{Var}_\gamma$ follows from $\sum_j\gamma_{ij}(\vartheta^0_j - \bar T_i) = 0$:
for each $i$, $\sum_j\gamma_{ij}\norm{\vartheta_i - \vartheta^0_j}^2
= \sum_j\gamma_{ij}\bigl(\norm{\vartheta_i - \bar T_i}^2 + \norm{\bar T_i - \vartheta^0_j}^2\bigr)$. For the pure
leak, $Z$ is non-increasing and $D^+Z \le -\sigma_mV_{\mathrm{bar}}(t)^2$, with $V_{\mathrm{bar}}$ computed from the
plan in use (the coupling inequality above holds for every optimal plan, vertex or not), so $Z$ converges and
$\int_0^\infty V_{\mathrm{bar}}^2\dd t \le Z(0)/\sigma_m < \infty$. $V_{\mathrm{bar}}$ itself is neither monotone nor
continuous --- it jumps at plan switches, upward as well as downward (numerics below) --- so the $\omega$-limit
statement rests on the stability of the LP instead. Let $\vartheta_\infty$ be an $\omega$-limit point and $F_\infty$
the set of optimal plans at $\vartheta_\infty$ (for a vertex-selection solution, its finitely many vertices). For a
fixed plan $\gamma$ the barycentres $\bar T^\gamma_i$ do not depend on $\vartheta$, so
$V_{\mathrm{bar}}^\gamma(\vartheta) := \bigl(\frac1m\sum\norm{\vartheta_i - \bar T^\gamma_i}^2_{\Ginv}\bigr)^{1/2}$ is
$1$-Lipschitz in $\vartheta$ for the norm $\bigl(\frac1m\sum\norm{\cdot}^2_{\Ginv}\bigr)^{1/2}$; and the optimal face
is upper semicontinuous --- the plan costs are continuous in $\vartheta$ and the polytope has finitely many vertices
--- so every plan optimal at a $\vartheta$ close enough to $\vartheta_\infty$ lies in $F_\infty$. Suppose no plan of
$F_\infty$ is centroidal; then $\delta := \min_{\gamma\in F_\infty}V_{\mathrm{bar}}^\gamma(\vartheta_\infty) > 0$ (a
minimum over a compact set), and for $\norm{\vartheta - \vartheta_\infty} < \delta/2$,
$V_{\mathrm{bar}}(\vartheta) \ge \delta/2$ whatever optimal plan is in use. The speed is bounded by
$\sigma_m\cdot\diam_{\Ginv}(\ThK)$ (Lemma P.1), so the solution, which enters the ball of radius $\delta/4$ around
$\vartheta_\infty$ at arbitrarily large times, either stays in the ball of radius $\delta/2$ from some time on or
spends a time $\ge\delta/(4\sigma_m\diam)$ in it on each of infinitely many visits; either way
$\int V_{\mathrm{bar}}^2\dd t = \infty$, a contradiction. Hence some plan of $F_\infty$ --- a vertex for
vertex-selection solutions --- is centroidal at $\vartheta_\infty$. (A two-way tie between vertex plans $\gamma$ and
$\rho$ is not a sliding mode of the vertex law unless $\bar T^\gamma = \bar T^\rho$: the rate of change of
$c_\rho - c_\gamma$ when following $\gamma$, plus that of $c_\gamma - c_\rho$ when following $\rho$, equals
$(2\sigma_m/m)\sum_i\norm{\bar T^\gamma_i - \bar T^\rho_i}^2_{\Ginv} \ge 0$, so at least one of the two plans stays
optimal when followed. Ties between three or more plans are not analysed.) The gradient identity is the envelope
theorem for the transportation LP at a configuration with a unique optimal plan:
$\partial_{\vartheta_i}\bigl[\tfrac12\sum\gamma_{ij}\norm{\vartheta_i - \vartheta^0_j}^2_{\Ginv}\bigr]
= \sum_j\gamma_{ij}\Ginv(\vartheta_i - \vartheta^0_j) = \frac1m\Ginv(\vartheta_i - \bar T_i)$, so
$-\sigma_m(\vartheta_i - \bar T_i) = -\sigma_m m\,G\,\partial_{\vartheta_i}\bigl[\tfrac12W^2\bigr]$. (Where the plan
is not unique, $\frac12W^2(\cdot,\nu)$ is not differentiable; the generalised velocities are then
$-\sigma_m m\,G\,\xi$ with $\xi$ ranging over the barycentric projections of the optimal plans --- the elements of the
(sub)differential of $\frac12W^2(\cdot,\nu)$ at an atomic $\mu$, cf.\ Ambrosio--Gigli--Savaré~\cite[\S10.4]{ags2008}.)

\smallskip
\textit{(ii)} Relabel so that $\frac1m\sum\norm{\vartheta_i - \vartheta^*_i}^2_{\Ginv} = \eta^2$. For any plan $\rho$
with row sums $1/m$, the plan cost $c_\rho(\vartheta) := \sum\rho_{ij}\norm{\vartheta_i - \vartheta^0_j}^2_{\Ginv}$
is the squared norm of $(\vartheta_i - \vartheta^0_j)$ in $L^2(\rho)$, so
$\abs{\sqrt{c_\rho(\vartheta)} - \sqrt{c_\rho(\vartheta^*)}} \le \eta$ by the triangle inequality (the labels of
$\vartheta$ and $\vartheta^*$ are matched). The barycentric part is dominated by the plan cost:
$C_{\mathrm{bar}} = \frac1m\sum\norm{\vartheta^*_i - \bar T_i}^2 \le c_\gamma(\vartheta^*)$ because
$\sum_j\gamma_{ij}\norm{\vartheta^*_i - \vartheta^0_j}^2 = \frac1m\norm{\vartheta^*_i - \bar T_i}^2
+ \sum_j\gamma_{ij}\norm{\bar T_i - \vartheta^0_j}^2$. Let $\gamma^*$ be optimal for $(\vartheta^*,\nu)$. Then
\[
  \sqrt{C_{\mathrm{bar}}} \le \sqrt{c_\gamma(\vartheta^*)} \le \sqrt{c_\gamma(\vartheta)} + \eta
  \le \sqrt{c_{\gamma^*}(\vartheta)} + \eta \le \sqrt{c_{\gamma^*}(\vartheta^*)} + 2\eta = b(t) + 2\eta.
\]
For the equality case: $c_\gamma(\vartheta^*) - b(t)^2 \le (b+2\eta)^2 - b^2 = 4\eta(b+\eta) < \Delta_{\mathrm{gap}}(t)$;
the vertex $\gamma$ (a vertex by the choice of the law) has cost at $\vartheta^*$ strictly below the next vertex cost,
so it is optimal at $\vartheta^*$; if $\gamma^*(t)$ is unique, $\gamma = \gamma^*(t)$ and
$C_{\mathrm{bar}} = \frac1m\sum\norm{\vartheta^*_i - \bar T^*_i}^2 = b_{\mathrm{bar}}(t)^2$. [For a generalised
solution whose plan at $t$ is not a vertex, the same chain applies to each vertex of the optimal face at $\vartheta$,
so every such vertex is optimal at $\vartheta^*$, and uniqueness of $\gamma^*(t)$ again gives $\gamma = \gamma^*(t)$.]

\smallskip
\textit{(iii)} Apply Theorem B with $y = \nabla g_t - \sigma_m\Ginv(\vartheta_i - \bar T_i)$ in (P). The extra term is
\[
  -\sigma_m\frac1m\sum\ip{\tilde\vartheta_i}{\tilde\vartheta_i + \vartheta^*_{\pi^*(i)} - \bar T_i}_{\Ginv}
  \le -\sigma_m\eta^2 + \sigma_m\eta\sqrt{C_{\mathrm{bar}}}.
\]
Under the gap condition $\sqrt{C_{\mathrm{bar}}} = b_{\mathrm{bar}}(t) \le b_{\mathrm{bar}}$ and
$\sigma_m\eta b_{\mathrm{bar}} \le (\sigma_m/2)(\eta^2 + b_{\mathrm{bar}}^2)$.

\smallskip
\textit{(iv)} is the proof of Theorem C(iv) with $b_{\mathrm{bar}}$ in place of $b$.

\smallskip
\textit{(v)} $b_{\mathrm{bar}}(t)^2 = W(\mu^*_t,\nu)^2 - \mathrm{Var}_{\gamma^*(t)} \le W(\mu^*_t,\nu)^2$, with
equality iff $\mathrm{Var}_{\gamma^*} = 0$, that is, iff every cell of $\gamma^*$ consists of a single atom of $\nu$
(coincident atoms counted once) --- the case $m_0 = m$, but also $m_0\mid m$ and priors with repeated atoms --- and
$b_{\mathrm{bar}} = 0$ iff $\vartheta^*_j = \bar T^*_j$ for all $j$ (at a tie, with the plan $\gamma^*(t)$ fixed by
the selection).
\end{proof}

\begin{rmk}{Remark}{C$'$.1}[reading]
The barycentric leak replaces ``pull each kernel to its prior kernel'' by ``pull each kernel to the centre of mass of
the prior it is responsible for''. Three consequences.
\begin{itemize}
  \item \emph{Transfer across sizes.} A network of $m$ kernels can be anchored to a learned measure of $m_0\ne m$
    kernels. The bias in the Lyapunov inequality is $b_{\mathrm{bar}}$, the distance of the ideal from the centroids
    of its cells, not the quantisation error $W(\mu^*,\nu)$ itself.
  \item \emph{Quantisation on the fly.} With $e\equiv0$ the law is Lloyd's algorithm for the equal-mass quantisation
    of $\nu$ in the $G^{-1}$ metric, in continuous time --- the continuous-time Lloyd descent of coverage
    control~\cite{S18}, with the equal-mass constraint of capacity-constrained quantisation~\cite{S19}; so the
    controller's leak and the quantisation theory of Corollary B.6 (Section~\ref{sec:program}, branch 9) are the same
    object. Like Lloyd's algorithm it converges to a critical point of the quantisation energy, not necessarily the
    global minimiser.
  \item \emph{Theorem F.} The leak clauses (i$'$)--(iii$'$) of Theorem F hold with $W(\mu_r,\nu)$ replaced by
    $V_{\mathrm{bar}}(r) \le W(\mu_r,\nu)$, since
    $\norm{\dot\vartheta_i}_{\Ginv} \le \norm{\nabla g_t}_G + \sigma_m\norm{\vartheta_i - \bar T_i}_{\Ginv}$ by
    Lemma P.1 and Minkowski.
\end{itemize}
\end{rmk}

\paragraph{Checked} (\code{tests/test_barycentric.py}). The plan reduces to the assignment for $m_0 = m$
($10^{-10}$) and the Pythagorean split $W^2 = V_{\mathrm{bar}}^2 + \mathrm{Var}_\gamma$ holds ($10^{-10}$); along the
pure leak $W^2$ is non-increasing at every step, the Dini bound $D^+Z \le -\sigma_mV_{\mathrm{bar}}^2$ holds at every
state (within 5\% of the first-order Euler quotient), the flow reaches $V_{\mathrm{bar}}^2 < 10^{-8}$, and the end
point is a fixed point of the equal-mass Lloyd iteration with the same $W^2$; $\sqrt{C_{\mathrm{bar}}} \le b + 2\eta$
on 90 random configurations, with $C_{\mathrm{bar}} = b_{\mathrm{bar}}^2$ at small $\eta$; the Lyapunov inequality
(iii) holds at every recorded state of a closed loop with the barycentric leak toward a 24-atom prior. The numbers are
in \code{results/barycentric_leak/table.md} (\code{exp_barycentric_leak.py}).

\paragraph{Numerics} (\code{results/barycentric_leak/table.md}).
\begin{itemize}
  \item \emph{Pure leak toward a 48-atom prior} ($m = 6$, lifted dimension 3, $G = \operatorname{diag}(2,1,0.5)$,
    $\sigma_m = 1$, three random starts). $W^2(\mu_t,\nu)$ is non-increasing at every one of the 3000 steps; the flow
    stops where the centroid defect is at the rounding level ($V_{\mathrm{bar}}^2 \le 2.2\cdot10^{-25}$) and
    $W^2 = \mathrm{Var}_\gamma$ ($1.164$, $1.213$, $1.203$ from $2.55$, $2.30$, $2.58$), after 2--6 plan switches, at
    which $V_{\mathrm{bar}}$ jumps (upward as well as downward). The end points are fixed points of the equal-mass
    Lloyd iteration; the best of 20 Lloyd runs from random starts reaches $1.171$, below two of the three end points
    and above the third: like Lloyd's algorithm, the leak finds a local minimiser of the quantisation energy, not
    always the global one. With $m_0 = m$ the same code reproduces the exact decay of Theorem C(i) to
    $5\cdot10^{-16}$.
  \item \emph{Closed loop} (representable plant with an $m^* = 4$ ideal, $m = 4$, $k_e = 4$, $\sigma_m = 0.5$, three
    seeds). Three leaks of the same gain: (A) the OT leak toward the ideal itself, (B) the barycentric leak toward a
    32-atom cloud of jittered copies of the ideal atoms, (C) the labelled leak toward a random 4-subset of the cloud.
    The network converges, in each case, to the configuration its leak anchors: (A) to the ideal ($W(\mu_T,\mu^*)$ and
    $\RMS\norm{e}$ below the printed precision); (B) to the equal-mass centroids of the cloud, at
    $W(\mu_T,\mu^*) = 0.167$, $0.097$, $0.157$, which equal the centroid defects $b_{\mathrm{bar}}$ of the ideal
    ($0.168$, $0.097$, $0.159$) although $W(\mu^*,\nu) = 0.45$, $0.33$, $0.40$; (C) to the subset, at
    $W(\mu_T,\mu^*) = 0.61$, $1.02$, $0.35$ against $W(\mu^*,\text{subset}) = 0.65$, $1.03$, $0.40$. The steady
    tracking RMS is $0.002$--$0.007$ for (B) and $0.03$--$0.04$ for (C). So the loop settles at the centroidal
    configuration of its prior, at distance $b_{\mathrm{bar}}$ from the ideal, as the fixed-point statements (i) and
    (v) describe; the Lyapunov bound (iv) does not apply to these runs (scope paragraph: the exchange bound on
    $\Delta_{\mathrm{gap}}$ and $4\eta(0)(b+\eta(0))$ are in the table). The comparison with (C) shows the cost of a
    prior of the right size at the wrong positions, not an advantage of barycentric over labelled leaks as such: a
    labelled leak toward the cloud's own equal-mass $m$-quantiser would have the same bias $b_{\mathrm{bar}}$ by (v).
\end{itemize}
